\section{Execution Times}

\subsection{Rust}

Obtaining the \textit{Worst-Case Execution Times} for every operation in the
system is a crucial step in order to perform response-time analysis. 

We decided to leverage the \textit{Data Watchpoint and Trace Unit (DWT)}
\cite[Chapter~9]{armlimitedARMCortexM42009}. \textit{DWT} supports multiple
functionalities such as: debugger hardware watchpoints, multiple comparators,
and cycle counters. For our purposes, the latter one was sufficient. 

In particular, the \textit{DWT} peripheral holds a cycle counter that can be
read from the \texttt{DWT\_CYCCNT} register, which in turn can be used to
implement a simple \textit{stopwatch} that can measure time spans. 

The chosen \textit{HAL} \cite{Stm32rsStm32f4xxhal2025} already provides a
stopwatch implementation that correctly handles the wrap around of the clock
counter. The typical system clock frequency of our board is \unit[168]{MHz},
which makes the period $T = \frac{1}{f} \simeq \unit[5.95]{ns}$. Since the
counter size is \unit[32]{bits}, it will wrap around every $2^{32} \cdot T
\simeq \unit[25]{s}$, thus it must be handled to prevent wrong readings. 

The stopwatch implementation provides an API that closely resembles the workings
of a manual stopwatch: it exports a \texttt{reset} function that simply clears
every stored time, and a \texttt{lap} method that stores the current cycle count
into a buffer passed during initialization. Lap times can be simply obtained by
subtracting successive cycle counts, having care of handling wrap around. 

Since this implementation is really basic, we have designed a higher-level
abstraction that allows us to obtain the WCET of selected code sections,
corresponding to the operations under analysis. 
Listing [\ref{lst:profiler-abstraction}] shows the \textit{trait} definition for
our profiler abstraction. The \texttt{lap} method gets called at the end of
every iteration (i.e. job execution), using the saved laps; then results can be
logged using the corresponding method. 

Every different task gets a tailored implementation of the profiler
\textit{trait}; they mainly differ on the implementation of the
\texttt{update\_wcet} method, to accommodate the different operations inside the
tasks. 

\begin{listing}
  \begin{minted}[breaklines]{rust}
    pub trait Profiler {
      fn reset(&mut self);
      fn lap(&mut self);
      fn lap_time(&self, lap: usize) -> Option<ProfilingDuration>;
      fn update_wcet(&mut self);
      fn log(&self);
  }
  \end{minted}
  \caption[Profiler abstraction code]{Profiler trait used to extract WCET}
  \label{lst:profiler-abstraction}
\end{listing}

The profiling instrumentation is gated behind \texttt{cargo} \textit{features}:
when the user desires to start the profiling of a certain operation it needs to
enable the related feature at compile-time. This ensures that the application is
completely devoid of profiling objects and instructions that are not needed
during normal operation, improving performance and code size. 

The runtime does not provide a \textit{task time} like Ada Ravenscar, hence, to
obtain the desired WCET, we had to manually raise the priority of the task under
analysis in order to prevent any interference that could invalidate our times.
This is not a very clean solution, but implementing a proper task execution time
would require non-trivial modifications to the RTIC runtime, which were
considered out of the scope of the work and too time-consuming. Nonetheless, the
WCET obtained remain valid. 

We have decided to take the maximum execution time among one-thousand
executions; when the correct profiling feature is enabled, the tasks are
configured to automatically log the WCET results accumulated by the profiler and
then panic to stop execution. 

\subsection{Ada}

Concerning the Ada system, we have re-utilized the instrumentation implemented
by J. Hu and A. Gobbo in their work \cite{huJiayihuSRT2024}, aimed to perform
response-time analysis over the \textit{Ada} system model, as shown in the
original paper \cite{burnsGuideUseAda2004}.

\subsection{WCET Times}

This section contains the obtained WCET values. Both implementations have been
compiled in \textit{release} mode (i.e. with all optimizations enabled), and run
on the board at maximum clock speed (\unit[168]{MHz}). 

We have observed some significant differences in execution times for certain
operations between the Ada and Rust system. We ascribe these gaps to differences
in I/O backends and runtime overheads, since RTIC heavily relies on hardware
features specific to ARM Cortex architecture, to compiler choices, and different
math libraries. 

\subsubsection{Activation Log Reader}
\begin{table}[H]
  \centering
  \begin{tabular*}{\linewidth}{@{\extracolsep{\fill}}lrr@{}}
    \toprule
    \textbf{Operation} & \textbf{Rust (\unit{ns})} & \textbf{Ada (\unit{ns})} \\
    \midrule
    Small Whetstone (workload 139) & 4346391 & 1018910 \\
    AL Read & 136 & 1267 \\
    ALR Read & 647 & 3256 \\
    Put Line & 8017 & 8222 \\
    \bottomrule
  \end{tabular*}
  \caption{WCET per operation — Activation Log Reader}
  \label{tab:wcet-activation-log-reader}
\end{table}

\subsubsection{External Event Server}
\begin{table}[H]
  \centering
  \begin{tabular*}{\linewidth}{@{\extracolsep{\fill}}lrr@{}}
    \toprule
    \textbf{Operation} & \textbf{Rust (\unit{ns})} & \textbf{Ada (\unit{ns})} \\
    \midrule
    AL Write & 583 & 1678 \\
    EES Write & 1071 & 3322 \\
    \bottomrule
  \end{tabular*}
  \caption{WCET per operation — External Event Server}
  \label{tab:wcet-external-event-server}
\end{table}

\subsubsection{On-Call Producer}
\begin{table}[H]
  \centering
  \begin{tabular*}{\linewidth}{@{\extracolsep{\fill}}lrr@{}}
    \toprule
    \textbf{Operation} & \textbf{Rust (\unit{ns})} & \textbf{Ada (\unit{ns})} \\
    \midrule
    RB Extract & 303 & 1361 \\
    Extract Workload & 814 & 4760 \\
    Small Whetstone (workload 278) & 8692782 & 2037820 \\
    Put Line & 7970 & 7406 \\
    \bottomrule
  \end{tabular*}
  \caption{WCET per operation — On-Call Producer}
  \label{tab:wcet-on-call-producer}
\end{table}

\subsubsection{Regular Producer}
\begin{table}[H]
  \centering
  \begin{tabular*}{\linewidth}{@{\extracolsep{\fill}}lrr@{}}
    \toprule
    \textbf{Operation} & \textbf{Rust (\unit{ns})} & \textbf{Ada (\unit{ns})} \\
    \midrule
    Small Whetstone (workload 756) & 23639364 & 5541700 \\
    Due Activation & 214 & 1344 \\
    RB Deposit & 2398 & 3583 \\
    OCP Activation & 2903 & 5283 \\
    Check Due & 208 & 1372 \\
    ALR Signal & 1982 & 6922 \\
    Cancel Deadline & 273 & 2835 \\
    RP Cancel Deadline & 8314 & -- \\
    \bottomrule
  \end{tabular*}
  \caption{WCET per operation — Regular Producer}
  \label{tab:wcet-regular-producer}
\end{table}

\subsubsection{Kernel/Runtime Overheads}
\begin{table}[H]
  \centering
  \begin{tabular*}{\linewidth}{@{\extracolsep{\fill}}lrr@{}}
    \toprule
    \textbf{Operation} & \textbf{Rust (\unit{ns})} & \textbf{Ada (\unit{ns})} \\
    \midrule
    Systick Overhead & 702 & 3906 \\
    ISR Switch & 2732 & 2850 \\
    Delay Until & 865 & 5420 \\
    Signal RTIC & 4160 & -- \\
    Task Semaphore & 4178 & -- \\
    Event Queue & 4101 & -- \\
    Spawn Overhead & 148 & -- \\
    Context-Switch & 166 & 3090 \\
    \bottomrule
  \end{tabular*}
  \caption{WCET — Kernel/Runtime Overheads}
  \label{tab:wcet-kernel-overheads}
\end{table}
